% Generated draft: requires mathematical and source review.
\documentclass[11pt]{article}
\usepackage[margin=25mm]{geometry}
\usepackage{fontspec}
\setmainfont{Latin Modern Roman}
\usepackage{amsmath,amssymb,mathtools,unicode-math}
\setmathfont{Latin Modern Math}
\usepackage{enumitem,needspace,xurl,hyperref}
\setlength{\parindent}{0pt}
\setlength{\parskip}{5pt}
\setlength{\emergencystretch}{3em}
\newcommand{\N}{\mathbb{N}}
\newcommand{\Z}{\mathbb{Z}}
\newcommand{\Q}{\mathbb{Q}}
\newcommand{\R}{\mathbb{R}}
\newcommand{\C}{\mathbb{C}}
\newcommand{\F}{\mathbb{F}}
\newcommand{\A}{\mathbb{A}}
\newcommand{\PP}{\mathbb{P}}
\newcommand{\E}{\mathbb{E}}
\newcommand{\eps}{\varepsilon}
\newcommand{\dd}{\,\mathrm{d}}
\begin{document}
\section*{UnsolvedMath ID: 800075}
\textbf{Source problem code:} catalogue entry
\subsection{Definitions and mathematical statement}
For a smooth volume form \(\eta\) on \(S^3\), write
\[
G_\eta=\operatorname{Diff}_0(S^3,\eta)
=\bigl\{f\in\operatorname{Diff}^{\infty}(S^3):f^*\eta=\eta\bigr\}_0,
\]
where the subscript denotes the identity component, interpreted here using the \(C^{\infty}\) topology. The source explicitly introduces a volume form only on \(S^2\), although its two questions concern \(S^3\). Thus \(\eta\) supplies the missing three-dimensional setup; the smoothness convention, topology, and intended quantification over volume forms require confirmation. The source's assertion of independence from the volume form is stated for \(S^2\), not explicitly for \(S^3\).

For a group \(G\), put \([a,b]=aba^{-1}b^{-1}\). The group is \emph{uniformly perfect} if
\[
\exists N\in\Z_{\geq 0}\ \forall g\in G\ \exists k\in\{0,\ldots,N\}\ \exists a_1,b_1,\ldots,a_k,b_k\in G:
\quad g=\prod_{i=1}^{k}[a_i,b_i],
\]
with the empty product equal to the identity. A real-valued \emph{quasimorphism} on \(G\) is a map \(q:G\to\R\) with finite defect
\[
D(q)=\sup_{g,h\in G}|q(gh)-q(g)-q(h)|<\infty.
\]
No continuity condition is imposed in this draft; the source does not specify one.

\textbf{Questions attributed to Burago, Ivanov, and Poltervoich, as named in the supplied source.}
\begin{enumerate}
\item Is \(G_\eta\) uniformly perfect?
\item Does \(G_\eta\) admit a non-trivial quasimorphism?
\end{enumerate}
The second question retains an unresolved terminology issue: the source does not define ``non-trivial.'' Two precise possible interpretations are existence of an unbounded quasimorphism, or existence of a quasimorphism that is not a bounded perturbation of a homomorphism. The latter asks for a quasimorphism \(q\) satisfying
\[
\forall \phi\in\operatorname{Hom}(G_\eta,\R),\qquad
\sup_{g\in G_\eta}|q(g)-\phi(g)|=\infty.
\]
These interpretations are recorded as alternatives, without identifying either as the verified intended formulation.

\textbf{Supplied background.} For a volume form \(\omega\) on \(S^2\), the source states that the isomorphism type of \(\operatorname{Diff}_0(S^2,\omega)\) is independent of \(\omega\). It attributes simplicity of this group to Banyaga, reference label \texttt{banyaga1978structure}, and the existence of many non-trivial quasimorphisms, with resulting failure of uniform perfectness, to Ghys and Gambaudo, reference label \texttt{gambaudo2004commutators}. These are supplied background assertions, not independently verified results in this draft. The registry labels the questions open; this is a supplied status label rather than an independently verified literature assessment.
\subsection{Short English statement}
The source asks whether the identity component of the volume-preserving diffeomorphism group of \(S^3\) has a uniform bound on the number of commutators needed to express any element, and whether it admits a non-trivial quasimorphism. The volume form on \(S^3\), the precise meaning of non-triviality, and regularity conventions need confirmation. This is a statement-only draft requiring human mathematical review.
\subsection{Sources}
\begin{itemize}
\item Local UnsolvedMath registry, supplied complete record: permanent ID \(800075\), exact source problem code \texttt{catalogue entry}, title ``Foliations and Diffeomorphism Groups --- Question 16.'' The record attributes the source list to Nariman and Yazdi, gives proposed year 2025, and supplies status \texttt{open}. Its creation and update timestamps are both \texttt{2026-10-16T14:00:00Z}.
\item Nariman and Yazdi, \emph{Problem List on Foliations and Diffeomorphism Groups} (2025), Question 16. Supplied source URL: \url{https://arxiv.org/abs/2503.18274}. The record reports access on 2026-10-16 and extraction by \texttt{arxiv\_latex}; these are supplied provenance statements, not a claim of access or independent verification for this draft. Supplied upstream revision: \texttt{d371ee917ac9647fbb06cd6fc529deec6f3dda82}.
\item Banyaga, source reference label \texttt{banyaga1978structure}: cited in the supplied statement for simplicity of the volume-preserving identity component on \(S^2\). Full bibliographic details are not supplied.
\item Ghys and Gambaudo, source reference label \texttt{gambaudo2004commutators}: cited in the supplied statement for non-trivial quasimorphisms on the corresponding \(S^2\) group and failure of uniform perfectness. Full bibliographic details are not supplied.
\item Burago, Ivanov, and Poltervoich, source reference label \texttt{burago2008conjugation}: credited by the supplied statement with posing the two \(S^3\) questions. The spelling ``Poltervoich'' is retained from the record; full bibliographic details are not supplied.
\end{itemize}
\end{document}
